% AUTO-GENERATED by scripts/exp_real_input_40_time_correlation_phase.py
\begin{table}[H]
\centering
\small
\setlength{\tabcolsep}{8pt}
\renewcommand{\arraystretch}{1.15}
\caption{真实输入 40 状态核在 $\chi$ 切片 $\theta(t)=(t,0,-t)$ 下的次主模态相位读出。记 $\lambda_1(t)$ 为 Perron--Frobenius 主特征值，$\lambda_2(t)$ 为其余谱中模最大的特征值之一，并令 $\rho_2(t)=|\lambda_2(t)|/\lambda_1(t)$，$\omega(t)=\arg(\lambda_2(t))$（主辐角）。若 $\omega(t)\ne 0$，则由 $(\lambda_2/\lambda_1)^n$ 的因子可定义内生振荡周期 $T_{\mathrm{osc}}(t)=2\pi/|\omega(t)|$。}
\label{tab:real-input-40-time-corr-phase}
\begin{tabular}{r r r r r r}
\toprule
 $t$ & $\lambda_1$ & $\Re\lambda_2$ & $\Im\lambda_2$ & $\rho_2$ & $T_{\mathrm{osc}}$\\
\midrule
-1.000 & 2.040720410275 & -1.217448524330 & 0 & 0.596577815462 & 2.000000\\
-0.500 & 2.287538746528 & -1.389014262764 & 0 & 0.607209064709 & 2.000000\\
0 & 2.618033988750 & -1.618033988750 & 0 & 0.618033988750 & 2.000000\\
0.500 & 3.059600073602 & -1.895186659885 & 0 & 0.619423001142 & 2.000000\\
1.000 & 3.646591323038 & -2.352028626673 & 0.182276700659 & 0.646927734108 & 2.050481\\
\bottomrule
\end{tabular}
\end{table}
